\documentclass[12pt]{article}

\newcommand{\duedate}{-----}
\newcommand{\assignment}{Exam 2 Summary}

% Change the following to your name and UNI.
\newcommand{\name}{Aksel Kretsinger-Walters, adk2164}
\newcommand{\email}{}

% NOTE: Defining collaborators is optional; to not list collaborators, comment out the
% line below. Maximum of two collaborators per problem set.
%\newcommand{\collaborators}{Vig Vigerton (\texttt{UNI}), Alice Bob (\texttt{UNI})}
% No collaborators on PS 0

\makeatletter
\def\input@path{{../}{../../}{../../../}} % add as many parents as you need
\makeatother
\input{pset_template_NNDL.tex} %% DO NOT CHANGE THIS LINE

\begin{document}
\psetheader %% DO NOT CHANGE THIS LINE

\section{RNN + Attention vs.\ Transformer Architectures}

\subsection{RNN + Attention (Seq2Seq)}
The classical sequence--to--sequence architecture consists of:
\begin{itemize}
		\item An \textbf{encoder RNN} (LSTM or GRU) that processes the input sequence sequentially:
    \[
						h_t^{\text{enc}} = \mathrm{RNN}(x_t,\, h_{t-1}^{\text{enc}}).
    \]
		\item A \textbf{decoder RNN} that autoregressively generates output tokens:
    \[
						s_t = \mathrm{RNN}(y_{t-1},\, s_{t-1}).
    \]
		\item \textbf{Additive or dot-product attention} over encoder states:
    \[
						e_{t,i} = f_{\mathrm{att}}(s_{t-1},\, h_i^{\text{enc}}), \qquad
						\alpha_{t,i} = \mathrm{softmax}_i(e_{t,i}),
    \]
    \[
						c_t = \sum_{i} \alpha_{t,i}\, h_i^{\text{enc}}.
    \]
		\item The decoder output distribution is conditioned on both the RNN state and context:
    \[
						p(y_t \mid y_{<t}, x) = \mathrm{softmax}(W[s_t, c_t]).
    \]
\end{itemize}

\paragraph{Key properties.}
\begin{itemize}
		\item Sequential computation in both encoder and decoder.
		\item Attention is \emph{cross-attention only}: decoder attends to encoder states.
		\item Long-range dependencies handled via recurrent states + attention.
		\item Parallelism limited: $O(n)$ sequential steps.
\end{itemize}

\subsection{Transformer}
The Transformer removes recurrence entirely and relies on:
\begin{itemize}
		\item \textbf{Self-attention} in encoder and decoder layers:
    \[
						\mathrm{SelfAttn}(Q,K,V)
						= \mathrm{softmax}\!\left(\frac{QK^\top}{\sqrt{d_k}}\right)V,
    \]
    where $Q=K=V$ are token representations of the same layer.
		\item \textbf{Cross-attention} in the decoder, where
    \[
						Q = H^{\text{dec}}, \quad K=V = H^{\text{enc}}.
    \]
		\item \textbf{Position-wise feed-forward networks} applied identically at each position.
		\item \textbf{Positional encodings} to inject order information.
\end{itemize}

\paragraph{Key properties.}
\begin{itemize}
		\item No recurrence: full parallelism across sequence length.
		\item Encoder uses stacked self-attention layers only.
		\item Decoder uses \emph{masked} self-attention + cross-attention.
		\item Complexity dominated by attention: $O(n^2)$ per layer.
\end{itemize}

\subsection{High-Level Comparison}
\begin{center}
		\begin{tabular}{lcc}
				\toprule
				& \textbf{RNN + Attention} & \textbf{Transformer} \\
				\midrule
				Core operation & Recurrence & Self-attention \\
				Parallelism & Low & High \\
				Long-range modeling & Via hidden state + attention & Direct via self-attention \\
				Attention type & Cross-attention only & Self-attn + cross-attn (decoder) \\
				Order information & Implicit in recurrence & Positional encodings \\
				Complexity & $O(n)$ sequential steps & $O(n^2)$ per layer \\
				\bottomrule
		\end{tabular}
\end{center}

\begin{figure}[h!]
		\centering
		\begin{tikzpicture}[
						node distance=1.2cm,
						box/.style={draw, rounded corners, minimum width=2.6cm, minimum height=0.9cm, align=center},
						arr/.style={->, thick}
				]

				% RNN + Attention (left)
				\node[box] (enc) {Encoder RNN};
				\node[box, below=of enc] (att) {Attention \\ (cross-attn)};
				\node[box, below=of att] (dec) {Decoder RNN};

				\draw[arr] (enc) -- (att);
				\draw[arr] (att) -- (dec);

				% Divider
				\node at (4.3,-1) {\Large $\Longleftrightarrow$};

				% Transformer (right)
				\node[box, right=4.8cm of enc] (encT) {Encoder Layer\\ (Self-Attn)};
				\node[box, below=of encT] (decT1) {Decoder Layer\\ (Masked Self-Attn)};
				\node[box, below=of decT1] (decT2) {Cross-Attn Layer};

				\draw[arr] (encT) -- (decT1);
				\draw[arr] (decT1) -- (decT2);

		\end{tikzpicture}
		\caption{Architecture comparison between RNN + Attention (left) and Transformer (right).}
\end{figure}

\end{document}
